% ============================================================
% Bakalarska praca - analyza komunikacnych protokolov (standalone)
% Datum benchmarku: 2026-03-23
% ============================================================

\documentclass[12pt,a4paper]{report}

\usepackage[T1]{fontenc}
\usepackage[utf8]{inputenc}
\usepackage[slovak]{babel}
\usepackage{lmodern}
\usepackage{geometry}
\geometry{margin=2.5cm}
\usepackage{amsmath}
\usepackage{hyperref}

\begin{document}

\chapter{Kvantitativna analyza prenosovej efektivity komunikacnych metod vo Fleet v2}
\label{chap:analyza-protokolov-fleet-v2}

\section{Vyskumny ciel a hodnotene varianty}

Cielom tejto kapitoly je kvantitativne porovnat transportne metody
pouzitelne pre protokol Fleet v2:
\begin{itemize}
  \item \texttt{protobuf-only} (unidirectional (unary) prenos bez streamingu),
  \item \texttt{gRPC protobuf},
  \item \texttt{REST/JSON},
  \item \texttt{JSON-RPC}.
\end{itemize}

Vyskumna otazka je formulovana nasledovne:
\textit{O kolko sa zmeni prenosova a vypoctova efektivita pri prechode
z binarneho protobuf prenosu na JSON metody v rovnakom workloade?}

\section{Metodika experimentu}

\subsection{Definicia meraneho scenara}

Meranie bolo realizovane pre jednu standardnu synchronizacnu sekvenciu
\texttt{PREPARE+REPORT}, ktora obsahuje 10 unidirectional sprav:
\texttt{Hello}, \texttt{GetAssignment}, \texttt{GetPolicy},
\texttt{AckPrepared}, \texttt{PublishReport} (request/response).

Pouzite boli tri profilove velkosti workloadu podla poctu eventov v reporte:
\begin{itemize}
  \item \textbf{small}: 10 eventov,
  \item \textbf{medium}: 50 eventov,
  \item \textbf{large}: 200 eventov.
\end{itemize}

\subsection{Vyhodnocovacie metriky}

Pre kazdu variantu bola spocitana velkost prenasanych dat v bytoch.
Percentualna zmena je definovana vztahom:
\begin{equation}
\Delta_{A \rightarrow B} [\%] = \frac{A - B}{B} \cdot 100
\end{equation}
kde kladna hodnota znamena, ze varianta $A$ je vacsia ako $B$.

Priemery v kapitole su uvadzane ako aritmeticky priemer percentualnych
hodnot cez profily small/medium/large (bez vahovania podla velkosti profilu).
Percenta aj nasobky su zaokruhlene na dve desatinne miesta.

\subsection{Definicia symbolov a skratiek}

Pre lepsiu citatelnost su vsetky premenne a skratky pouzite v kapitole
zhrnute na jednom mieste:
\begin{itemize}
  \item $A$, $B$ [B]: porovnavane hodnoty objemu dat (v bytoch) pre dve metody,
  \item $\Delta_{A \rightarrow B}$ [\%]: relativna zmena hodnoty $A$ voci $B$,
  \item $r_{A<B}$ [\%]: percentualne skratenie casu prenosu metody $A$ voci $B$,
  \item $s_{A/B}$ [x]: nasobok rychlosti metody $A$ voci $B$ pri rovnakom bitrate,
  \item $t$ [s]: cas prenosu; pri rovnakom bitrate plati $t \propto \text{bytes}$,
  \item $N_{\text{msg}}$ [-]: pocet sprav v jednej sekvencii PREPARE+REPORT (tu 10),
  \item $N_{\text{rpc}}$ [-]: pocet unidirectional RPC volani v sekvencii (tu 5),
  \item $H_{\text{meta}}$ [B/RPC]: odhad velkosti HTTP/2 metadata na jedno RPC,
  \item $O_{\text{gRPC/H2}}$ [B]: modelovany overhead gRPC/HTTP2 na celu sekvenciu,
  \item $k$ [-]: pocet behov (opakovani) celej sekvencie PREPARE+REPORT,
  \item $\text{bytes}_{\text{proto-only}}$ [B]: celkovy objem dat
  pre \texttt{protobuf-only} v danom profile alebo scenari,
  \item $\Delta_{\text{gRPC-proto}}(k)$ [B]: rozdiel prenesenych dat
  (\texttt{gRPC wire} -- \texttt{protobuf-only}) pri $k$ behoch,
  \item $O_{\text{setup}}$ [B]: fixny overhead na zaciatku scenara
  (napr. inicializacia spojenia/kanala),
  \item $O_{\text{cycle}}$ [B/cycle]: prirastkovy overhead na jeden dalsi beh sekvencie.
\end{itemize}

Pouzite skratky:
\begin{itemize}
  \item RPC: Remote Procedure Call,
  \item RTT: Round-Trip Time,
  \item TLS: Transport Layer Security,
  \item HPACK: kompresia HTTP/2 hlaviciek,
  \item CPU: Central Processing Unit.
\end{itemize}

\subsection{Ukazkove vypocty}

V celej kapitole sa pouzivaju tri zakladne vztahy:
\begin{enumerate}
  \item \textbf{Relativna zmena velkosti} (so znamienkom):
  \begin{equation}
  \Delta_{A \rightarrow B} [\%] = \frac{A - B}{B} \cdot 100
  \end{equation}
  \item \textbf{Skratenie casu prenosu} (alebo zrychlenie):
  \begin{equation}
  r_{A<B} [\%] = \left(1 - \frac{A}{B}\right) \cdot 100
  \end{equation}
  \item \textbf{Nasobok rychlosti} (pri rovnakom bitrate):
  \begin{equation}
  s_{A/B} [x] = \frac{B}{A}
  \end{equation}
\end{enumerate}

\paragraph{Priklad 1: \texttt{protobuf-only} vs \texttt{REST/JSON} (medium).}
Z tabulky plati:
\texttt{protobuf-only}=10909~B a \texttt{REST/JSON}=23346~B.
\begin{equation}
\Delta = \frac{10909 - 23346}{23346} \cdot 100 = -53.27\%
\end{equation}
Znamienko minus znamena, ze \texttt{protobuf-only} je mensi.
Ak hodnotu interpretujeme ako skratenie casu prenosu:
\begin{equation}
r = \left(1-\frac{10909}{23346}\right)\cdot100 = 53.27\%
\end{equation}
Nasobok rychlosti:
\begin{equation}
s = \frac{23346}{10909} = 2.14x
\end{equation}
Teda \texttt{protobuf-only} je v profile medium o 53.27\% rychlejsi
na prenos (priblizne 2.14x) ako \texttt{REST/JSON}.

\paragraph{Priklad 2: overhead \texttt{gRPC/HTTP2} voci \texttt{protobuf-only} (small).}
Z tabulky plati:
\texttt{gRPC}=5384~B a \texttt{protobuf-only}=4509~B.
\begin{equation}
\Delta = \frac{5384 - 4509}{4509}\cdot100 = 19.41\%
\end{equation}
Ide o modelovany overhead \texttt{gRPC/HTTP2} voci \texttt{protobuf-only}.

\paragraph{Priklad 3: CPU zrychlenie protobuf serialize vs JSON.}
Z CPU tabulky plati:
\texttt{protobuf}=0.1438~s a \texttt{JSON}=1.1904~s.
\begin{equation}
r = \left(1-\frac{0.1438}{1.1904}\right)\cdot100 = 87.92\%
\end{equation}
To zodpoveda tvrdeniu, ze protobuf serializacia je o 87.92\% rychlejsia.

\subsection{Postup realizacie merania}

Aby bola metodika reprodukovatelna a auditovatelna, experiment bol realizovany
v tychto krokoch:
\begin{enumerate}
  \item Benchmark skript bol spusteny v kontajneri
  \texttt{monad-fleet-service}
  (\texttt{docker compose exec -T monad-fleet-service python ...}),
  teda v rovnakom prostredi, v akom bezi sluzba Fleet.
  \item Skript programovo vytvoril schema-validne spravy podla
  \texttt{fleet\_gateway\_v2.proto}:
  \texttt{Hello}, \texttt{GetAssignment}, \texttt{GetPolicy},
  \texttt{AckPrepared}, \texttt{PublishReport} (request/response).
  \item Profil workloadu sa menil poctom eventov v sprave
  \texttt{PublishReport}: 10 (small), 50 (medium), 200 (large).
  \item Pre kazdu spravu sa vypocitala velkost payloadu:
  \begin{itemize}
    \item protobuf bytes cez \texttt{SerializeToString()},
    \item JSON bytes cez \texttt{MessageToDict()} a
    \texttt{json.dumps(..., sort\_keys=True, separators=(',', ':'))}.
  \end{itemize}
  \item Transportne varianty boli modelovane nasledovne:
  \begin{itemize}
    \item \texttt{protobuf-only}: suma protobuf bytes,
    \item \texttt{gRPC protobuf (HTTP/2 model)}:
    \begin{equation}
    O_{\text{gRPC/H2}} =
    N_{\text{msg}}\cdot5 +
    N_{\text{rpc}}\cdot5\cdot9 +
    N_{\text{rpc}}\cdot H_{\text{meta}}
    \end{equation}
    kde $N_{\text{msg}}=10$, $N_{\text{rpc}}=5$ a
    $H_{\text{meta}}=120$~B na unidirectional RPC (modelovany odhad
    komprimovanych HTTP/2 metadata pre request/response/trailers),
    \item \texttt{REST/JSON}: suma JSON bytes,
    \item \texttt{JSON-RPC}: JSON-RPC envelope
    (\texttt{jsonrpc}, \texttt{id}, \texttt{method}, \texttt{params/result})
    + JSON payload.
  \end{itemize}
  \item CPU cast bola merana pre profil medium:
  5 behov x 8000 iteracii na kazdu operaciu
  (\texttt{protobuf serialize/parse}, \texttt{json serialize/parse},
  \texttt{json-rpc serialize/parse}).
  \item Vysledne hodnoty (payload aj CPU) boli ulozene do suboru
  \texttt{output/research/protocol\_transport\_benchmark\_2026-03-24\_http2\_modeled.json}
  (modelovana cast) a do suboru
  \texttt{output/research/unidirectional\_grpc\_wire\_benchmark\_2026-03-24.json}
  (realne unidirectional wire meranie), z ktorych su odvodene tabulky v tejto kapitole.
\end{enumerate}

\subsection{Ako bol urceny HTTP/2 overhead}

Pouzity overhead pre \texttt{gRPC/HTTP2} nebol zvoleny nahodne, ale bol
zostaveny z troch casti:
\begin{itemize}
  \item \textbf{gRPC message frame:} 5~B na spravu (1~B compression flag +
  4~B dlzka payloadu),
  \item \textbf{HTTP/2 frame header:} 9~B na frame; pre unidirectional RPC bol
  modelovany tok 5 framov
  (request HEADERS, request DATA, response HEADERS, response DATA, trailers),
  \item \textbf{HTTP/2 metadata:} odhad $H_{\text{meta}}=120$~B na unidirectional RPC
  po HPACK kompresii (request/response/trailers metadata).
\end{itemize}

Pre sekvenciu PREPARE+REPORT (10 sprav, 5 unidirectional RPC) to dava:
\begin{equation}
O_{\text{gRPC/H2}} = 10\cdot5 + 5\cdot5\cdot9 + 5\cdot120 = 875~\text{B}
\end{equation}

Keze presna metadata velkost zavisi od implementacie a stavu HPACK tabuliek,
bola doplnena citlivostna analyza:
\begin{table}[htbp]
\centering
\caption{Citlivost overheadu na predpokladanu metadata velkost}
\label{tab:grpc-http2-metadata-sensitivity}
\begin{tabular}{|r|r|r|r|r|}
\hline
$H_{\text{meta}}$ [B/RPC] & $O_{\text{gRPC/H2}}$ [B] & small [\%] & medium [\%] & large [\%] \\
\hline
60  & 575  & 12.75\% & 5.27\% & 1.65\% \\
120 & 875  & 19.41\% & 8.02\% & 2.51\% \\
180 & 1175 & 26.06\% & 10.77\% & 3.37\% \\
\hline
\end{tabular}
\end{table}

\subsection{Limity metodiky}

Tento experiment predstavuje payload/codec benchmark, nie packet-level
sietove meranie. To znamena, ze:
\begin{itemize}
  \item meria sa velkost tela sprav a CPU cena kodovania/dekodovania,
  \item nemeraju sa realne RTT, packet loss, TCP retransmisie, TLS handshake,
  HTTP/2 connection management ani spravanie siete pod zatazou,
  \item pri \texttt{gRPC} je HTTP/2 rezia modelovana odhadom
  $H_{\text{meta}}=120$~B na unidirectional RPC; nejde o packet-level meranie
  skutocnej siete.
\end{itemize}

Pre \texttt{protobuf-only} baseline sa v tejto kapitole predpoklada
priamy binarny prenos protobuf payloadu bez gRPC/HTTP2 obalenia.
To znamena, ze HTTP/2 connection management sa do \texttt{protobuf-only}
vysledkov nezapocitava. Ak by bol protobuf preneseny cez HTTP/2 kanal,
treba k nemu analogicky priratat aj HTTP/2/TLS rezijnu cast.

Vysledky su preto urcene predovsetkym na porovnanie relativnej efektivity
formatov a transportneho obalenia, nie na absolutnu predikciu end-to-end
latencie v produkcnej sieti.

Pre produkcnu predikciu je vhodne doplnit samostatne packet-level meranie:
\begin{itemize}
  \item connection reuse vs reconnect (vplyv TCP/TLS/HTTP2 setupu),
  \item paralelne klienty a vyssia zataz (head-of-line/blocking efekty),
  \item sietove podmienky (RTT, packet loss, retransmisie, jitter).
\end{itemize}

\subsection{Vplyv typu komunikacneho toku}

Aktualny Fleet v2 tok je unidirectional (v gRPC terminologii unary),
preto benchmark cieli na realny rezim aktualneho nasadenia.
Ak by bol tok navrhnuty ako dlhodoby bidirectional stream, vysledky by sa
pravdepodobne zmenili, pretoze by sa inak rozlozila transportna rezia
a inac by sa uplatnilo riadenie toku.

\section{Vysledky merania prenosoveho objemu}

\begin{table}[htbp]
\centering
\caption{Celkovy objem prenesenych dat na jednu sekvenciu PREPARE+REPORT}
\label{tab:transport-bytes-profile}
\begin{tabular}{|l|r|r|r|r|}
\hline
Profil & protobuf-only [B] & gRPC protobuf (HTTP/2 model) [B] & REST/JSON [B] & JSON-RPC [B] \\
\hline
small  & 4509  & 5384  & 9796  & 10313 \\
medium & 10909 & 11784 & 23346 & 23863 \\
large  & 34910 & 35785 & 74159 & 74676 \\
\hline
\end{tabular}
\end{table}

Z nameranych dat vyplyva:
\begin{itemize}
  \item \texttt{gRPC protobuf} je voci \texttt{protobuf-only} vacsi
  v priemere o \textbf{9.98\%},
  \item \texttt{protobuf-only} je v priemere o \textbf{53.39\% mensi}
  ako \texttt{REST/JSON},
  \item \texttt{protobuf-only} je v priemere o \textbf{54.60\% mensi}
  ako \texttt{JSON-RPC},
  \item \texttt{JSON-RPC} je v priemere o \textbf{2.73\% vacsi}
  ako \texttt{REST/JSON}.
\end{itemize}

\section{Interpretacia vysledkov z hladiska casu prenosu}

Pri rovnakom bitrate je cas prenosu priamo umerny velkosti payloadu
($t \propto \text{bytes}$). Zo zmeranych bytov je preto mozne odvodit,
o kolko je jednotliva metoda rychlejsia alebo pomalsia na prenos.

\subsection{Porovnanie rychlosti voci \texttt{REST/JSON}}

Tabulka~\ref{tab:transfer-faster-vs-rest} uvadza skratenie casu prenosu
voci \texttt{REST/JSON}. Kladna hodnota znamena kratsi cas prenosu
(rychlejsie), zaporna hodnota znamena dlhsi cas (pomalsie).

\begin{table}[htbp]
\centering
\caption{Skratenie casu prenosu voci REST/JSON [\%]}
\label{tab:transfer-faster-vs-rest}
\begin{tabular}{|l|r|r|r|}
\hline
Profil & protobuf-only & gRPC protobuf (HTTP/2 model) & JSON-RPC \\
\hline
small  & +53.97\% & +45.04\% & -5.28\% \\
medium & +53.27\% & +49.52\% & -2.21\% \\
large  & +52.93\% & +51.75\% & -0.70\% \\
\hline
\end{tabular}
\end{table}

Ekvivalentne vyjadrenie v nasobku rychlosti (voci \texttt{REST/JSON}):
\begin{itemize}
  \item \texttt{protobuf-only}: \textbf{2.12x az 2.17x} rychlejsi prenos,
  \item \texttt{gRPC protobuf}: \textbf{1.82x az 2.07x} rychlejsi prenos,
  \item \texttt{JSON-RPC}: \textbf{0.95x az 0.99x}, teda mierne pomalsi
  prenos ako \texttt{REST/JSON}.
\end{itemize}

\subsection{Priame porovnanie \texttt{gRPC protobuf} vs \texttt{JSON-RPC}}

Pre priamu odpoved na volbu medzi \texttt{gRPC protobuf} a \texttt{JSON-RPC}
bol vypocitany pomer:
\begin{equation}
\text{rychlejsie}_{\text{gRPC vs JSON-RPC}}[\%] =
\left(1-\frac{\text{bytes}_{\text{gRPC}}}{\text{bytes}_{\text{JSON-RPC}}}\right)\cdot100
\end{equation}

\begin{table}[htbp]
\centering
\caption{Skratenie casu prenosu gRPC protobuf voci JSON-RPC [\%]}
\label{tab:grpc-vs-jsonrpc}
\begin{tabular}{|l|r|r|r|}
\hline
Profil & gRPC protobuf [B] & JSON-RPC [B] & gRPC rychlejsie [\%] \\
\hline
small  & 5384  & 10313 & 47.79\% \\
medium & 11784 & 23863 & 50.62\% \\
large  & 35785 & 74676 & 52.08\% \\
\hline
\end{tabular}
\end{table}

Z toho vyplyva, ze \texttt{gRPC protobuf} je v tomto workloade
objemovo aj casovo (pri rovnakom bitrate) priblizne \textbf{1.91x az 2.09x}
rychlejsi ako \texttt{JSON-RPC}.

\subsection{Preco je rozdiel medzi \texttt{protobuf-only} a \texttt{gRPC} vacsi}

Rozdiel medzi \texttt{protobuf-only} a \texttt{gRPC protobuf} je po zapocitani
HTTP/2 rezie vyraznejsi:
\begin{itemize}
  \item \texttt{protobuf-only} je rychlejsi o \textbf{2.51\% az 19.41\%}
  (podla velkosti profilu),
  \item priemerna hodnota je \textbf{9.98\%}.
\end{itemize}

Rozdiel vyplyva z modelu \texttt{gRPC/HTTP2}, ktory obsahuje:
\begin{equation}
O_{\text{gRPC/H2}} =
N_{\text{msg}}\cdot5 +
N_{\text{rpc}}\cdot5\cdot9 +
N_{\text{rpc}}\cdot H_{\text{meta}}
\end{equation}
\begin{equation}
= 10\cdot5 + 5\cdot5\cdot9 + 5\cdot120 = 875~\text{B}
\end{equation}
\begin{equation}
\Delta_{\text{gRPC vs proto}} [\%] =
\frac{875}{\text{bytes}_{\text{proto-only}}}\cdot 100
\end{equation}

Po dosadeni:
\begin{itemize}
  \item small: $\frac{875}{4509}\cdot100 = 19.41\%$,
  \item medium: $\frac{875}{10909}\cdot100 = 8.02\%$,
  \item large: $\frac{875}{34910}\cdot100 = 2.51\%$.
\end{itemize}

Absolutny rozdiel je v modeli rovnaky (875~B), ale relativny podiel klesa
s rastucou velkostou payloadu.

\subsection{Preco pri large profile nevznika velky rozdiel}

Pri large profile sa casto objavi dojem, ze by \texttt{gRPC} malo byt
vyrazne pomalsie, lebo nesie HTTP/2. Dolezite je vsak oddelit
\textbf{pocet RPC volani} od \textbf{velkosti payloadu}:
\begin{itemize}
  \item sekvencia PREPARE+REPORT ma stale 5 unidirectional RPC volani,
  \item \texttt{PublishReport} je jedno RPC, aj ked obsahuje 200 eventov,
  \item vyssi pocet eventov teda zvacsuje hlavne telo spravy, nie pocet
  HTTP request/response parov.
\end{itemize}

Preto fixna HTTP/2 rezia tvori pri large profile len malu cast celku
(v realnom wire merani pri $k=1000$ je rozdiel 2.12\%).
Ak sa v niektorej sade behov objavi opacne poradie v latencii,
ide typicky o vplyv serverovej biznis logiky, scheduler jitteru,
GC a inych runtime efektov, nie o zmenu efektivity samotneho codec-u.

\subsection{Porovnanie pri roznych poctoch behov (realne unidirectional meranie)}

Pre tuto cast bol vykonany samostatny realny unidirectional benchmark na wire urovni:
\begin{itemize}
  \item serverova strana bola realizovana ako lightweight
  \texttt{mock FleetManager} (rovnake \texttt{fleet.v2} protobuf spravy),
  \item \texttt{protobuf-only} je suma serializovanych request/response payloadov [B],
  \item \texttt{gRPC} je merane ako realne TCP payload byty [B] cez counting proxy
  medzi unidirectional klientom a serverom,
  \item artefakt merania:
  \texttt{output/research/unidirectional\_grpc\_wire\_benchmark\_2026-03-24.json}.
\end{itemize}

Z nameranych dat vyplyva, ze rozdiel je dobre aproximovany vztahom:
\begin{equation}
\Delta_{\text{gRPC-proto}}(k) [\text{B}] \approx O_{\text{setup}} + k \cdot O_{\text{cycle}}
\end{equation}
kde:
\begin{itemize}
  \item $k$ je pocet behov (opakovani) celej sekvencie PREPARE+REPORT,
  \item $O_{\text{setup}}$ je fixny overhead na zaciatku scenara,
  \item $O_{\text{cycle}}$ je prirastkovy overhead na jeden beh sekvencie.
\end{itemize}

Pre tento benchmark vyslo:
\begin{equation}
O_{\text{setup}} \approx 700~\text{B}, \quad O_{\text{cycle}} \approx 605~\text{B}
\end{equation}

Teda absolutny rozdiel rastie priblizne linearne s poctom behov $k$, ale
percentualny rozdiel klesa s rastom workloadu (amortizacia setup overheadu).

Pre medium profil:
\begin{table}[htbp]
\centering
\caption{gRPC vs protobuf-only pri roznych poctoch behov (medium, realne meranie)}
\label{tab:grpc-proto-runs-medium}
\begin{tabular}{|r|r|r|r|r|}
\hline
Pocet behov $k$ & protobuf-only [B] & gRPC wire [B] & Rozdiel [B] & Rozdiel [\%] \\
\hline
1    & 7400    & 8683    & 1283    & 17.34\% \\
5    & 37000   & 40703   & 3703    & 10.01\% \\
10   & 74000   & 80775   & 6775    & 9.16\% \\
50   & 370000  & 400996  & 30996   & 8.38\% \\
100  & 740000  & 801246  & 61246   & 8.28\% \\
1000 & 7400000 & 8006146 & 606146  & 8.19\% \\
\hline
\end{tabular}
\end{table}

Pre uplnost, celkovy prenos dat pri $k=1000$ behoch:
\begin{table}[htbp]
\centering
\caption{Prenesene data pri 1000 behoch PREPARE+REPORT (realne unidirectional meranie)}
\label{tab:transfer-1000-runs-all-profiles}
\begin{tabular}{|l|r|r|r|r|}
\hline
Profil & protobuf-only [B] & gRPC wire [B] & Rozdiel [B] & Rozdiel [\%] \\
\hline
small  & 1858000  & 2464120  & 606120 & 32.62\% \\
medium & 7400000  & 8006146  & 606146 & 8.19\% \\
large  & 28558000 & 29164099 & 606099 & 2.12\% \\
\hline
\end{tabular}
\end{table}

\section{Vysledky CPU analyzy}

CPU benchmark bol vykonany pre profil medium, 5 behov, 8000 iteracii na beh.

\subsection{Rozsah merania CPU}

Meranie hodnoti iba serializaciu a parsovanie payloadu
na aplikacnej vrstve (Python):
\begin{itemize}
  \item \texttt{protobuf}: \texttt{SerializeToString()} a \texttt{ParseFromString()},
  \item \texttt{JSON}/\texttt{JSON-RPC}: \texttt{json.dumps()} a \texttt{json.loads()}.
\end{itemize}

Benchmark \textbf{nemeria} sietovu latenciu, TLS handshake, kernel/network stack
ani retry logiku. Cielom je izolovat cistu CPU cenu kodovania dat.

\begin{table}[htbp]
\centering
\caption{Priemerne casy serializacie a parsovania}
\label{tab:cpu-serialize-parse}
\begin{tabular}{|l|r|r|r|}
\hline
Operacia & protobuf [s] & JSON [s] & JSON-RPC [s] \\
\hline
Serialize & 0.1438 & 1.1904 & 1.1827 \\
Parse     & 0.1355 & 0.7879 & 0.8300 \\
\hline
\end{tabular}
\end{table}

Odvodene percentualne zrychlenie protobuf:
\begin{itemize}
  \item serialize: \textbf{87.92\%} rychlejsie ako JSON,
  \item parse: \textbf{82.81\%} rychlejsie ako JSON,
  \item serialize: \textbf{87.84\%} rychlejsie ako JSON-RPC,
  \item parse: \textbf{83.68\%} rychlejsie ako JSON-RPC.
\end{itemize}

Ekvivalentne to znamena, ze JSON encode je priblizne \textbf{8.28x} pomalsi
a JSON parse priblizne \textbf{5.82x} pomalsi ako protobuf.

\subsection{Vztah CPU vysledkov k rozdielu \texttt{protobuf-only} vs \texttt{gRPC}}

Pri interpretacii je potrebne oddelit codec vrstvu od transportnej vrstvy:
\begin{itemize}
  \item \texttt{protobuf-only} aj \texttt{gRPC protobuf} pouzivaju rovnaky
  binarny protobuf codec, preto je CPU cena serializacie a parsovania
  prakticky rovnaka,
  \item rozdiel medzi nimi v tejto kapitole sposobuje najma transportne
  obalenie HTTP/2 (framing + metadata model), nie odlisny format payloadu.
\end{itemize}

Z tohto dovodu je rozdiel medzi \texttt{protobuf-only} a \texttt{gRPC}
v CPU rovine maly (rovnaky codec), ale v prenosovej rovine uz meratelny
v dosledku HTTP/2 rezie; rozdiel oproti JSON metodam vsak zostava vyrazny.

\section{Diskusia}

Vysledky je mozne zhrnut do troch hlavnych bodov:
\begin{enumerate}
  \item Pre objemovo narocnu telemetriu je binarny protobuf prenos
  vyrazne efektivnejsi ako JSON varianty.
  \item Rozdiel medzi \texttt{gRPC protobuf} a \texttt{protobuf-only}
  je sposobeny najma HTTP/2 reziou; po jej zapocitani je v benchmarku
  \texttt{protobuf-only} rychlejsi o 2.51\% az 19.41\%.
  \item \texttt{JSON-RPC} pridava oproti \texttt{REST/JSON}
  metadatovy envelope, co zvysuje celkovy objem prenesenych dat.
\end{enumerate}

\subsection{Predpokladany vplyv bidirectional toku}

Uvedene vysledky platia pre sekvenciu samostatnych unidirectional volani.
Pri bidirectional rezime by sa mohli percenta odlisovat z tychto dovodov:
\begin{itemize}
  \item \textbf{Amortizacia overheadu:} pri jednom dlhodobom streame sa
  fixna rezia rozlozi medzi vacsi pocet sprav.
  \item \textbf{Iny pomer payload a riadiacej signalizacie:} pri castych
  malych spravach moze byt relativny podiel metadat odlisny od unidirectional modelu.
  \item \textbf{Flow-control a backpressure:} streamovacie protokoly menia
  spravanie throughputu pri burstoch a pri pomalsom prijimaci.
  \item \textbf{Interakcny model:} bidirectional tok vie znizit latenciu
  priebeznych eventov, ale moze mat iny profil vyuzitia CPU a siete.
\end{itemize}

Prakticky dosledok je, ze poradie metod
(binarne protobuf varianty vs JSON varianty) pravdepodobne zostane rovnake,
avsak konkretne percentualne hodnoty sa mozu pri bidirectional scenari posunut.

\section{Zaver}

Pre Fleet v2 v aktualnom unidirectional rezime je najefektivnejsi
\texttt{protobuf-only} prenos. V pripade poziadavky na standardnu RPC vrstvu
predstavuje \texttt{gRPC protobuf} stale vyrazne efektivnejsiu volbu
nez JSON metody, avsak s meratelnou HTTP/2 reziou voci \texttt{protobuf-only}.

Rozdiel medzi \texttt{protobuf-only} a \texttt{gRPC} je sposobeny tym,
ze obe varianty pouzivaju rovnaky binarny protobuf codec, ale \texttt{gRPC}
zaroven vyuziva HTTP/2 transport. V pouzitom modeli to predstavuje overhead
875~B na celu sekvenciu PREPARE+REPORT (framing + metadata odhad),
t.~j. 2.51\% az 19.41\% podla profilu.

Realne unidirectional wire meranie navyse ukazalo, ze pri viacnasobnych behoch
sa overhead sprava ako
$\Delta_{\text{gRPC-proto}}(k) \approx O_{\text{setup}} + k\cdot O_{\text{cycle}}$,
kde v tomto experimente vyslo priblizne
$O_{\text{setup}}\approx700$~B a $O_{\text{cycle}}\approx605$~B.
Pri $k=1000$ bolo namerane:
small 32.62\%, medium 8.19\%, large 2.12\%.

V realnom nasadeni sa k tomu priratava dalsia gRPC/HTTP/2 rezia
(najma dynamika headers/trailers, HPACK, TLS a riadenie streamu), preto sa
konkretne percenta mozu odlisovat od modelovanych hodnot.

Pri prechode na bidirectional tok by sa konkretne percenta mohli zmenit,
najma v dosledku odlisnej amortizacie overheadu a dynamiky riadenia toku.
Z pohladu architektonickeho rozhodovania vsak benchmark konzistentne ukazuje,
ze hlavny rozdiel vznikne medzi binarnymi protobuf variantmi a JSON metodami,
nie medzi \texttt{protobuf-only} a \texttt{gRPC}.

\end{document}
